\naslov{Polinomska interpolacija}

\podnaslov{Lagrangeova oblika}

Problem z interpolacijo v standardni bazi je, da je Vandermondova matrika zelo
občutljiva.

\begin{izrek}
  Za paroma različne točke $x_0, \ldots, x_n$ in vrednosti $y_0, \ldots, y_n$
  obstaja natanko en polinom stopnje $\le n$, da je $p(x_i) = y_i$.
\end{izrek}

\begin{proof}
  Obstoj dokažemo s konstrukcijo, kjer uporabimo Lagrangeove bazne polinome
  \[
	l_{n,i}(x_j) =
	\begin{cases}
	  1 & i = j \\
	  0 & i \ne j
	\end{cases}.
  \]
  Če take polinome lahko konstruiramo (kjer je $l_{n.i}$ stopnje $\le n$), bo
  veljalo
  \[
	p = \sum_i y_i l_{n,i}.
  \]
  Če definiramo
  \[
	l_{n,i}(x)
	= \frac{(x-x_0) \ldots (x - x_{i-1}) (x - x_{i+1}) \ldots (x - x_n)}{(x_i -
	  x_0) \ldots (x_i - x_{i-1}) (x_i - x_{i+1}) \ldots (x_i - x_n)},
  \]
  potem $l_{0,n}, \ldots, l_{n,n}$ sestavljajo bazo za polinome stopnje $\le n$.

  Dokazati moramo še enoličnost:
  denimo, da je $\tilde{p}$ tudi polinom stopnje $\le n$, ki zadošča
  $\tilde{p}(x_i) = y_i$.
  Tedaj je $q = p - \tilde{p}$ polinom stopnje $\le n$, za katerega velja
  $q(x_i) = 0$.
  Polinom stopnje $\le n$, ki ima vsaj $n+1$ ničel, je enak $q = 0$.
\end{proof}

\vprasanje{Kako interpoliraš polinom z Lagrangeovo bazo? Dokaži, da je
  interpolacija enolična.}

% LocalWords:  Vandermondova
